\SourceLeaf{222}%
\bigskip\Needspace{10\baselineskip}
\pdfbookmark[1]{Fortgesetzte Nachrichten. Pallas Olbersiana. Juli 1802}{pallas180207}
\begin{center}
{\small FORTGESETZTE NACHRICHTEN ÜBER DEN NEUEN HAUPT-PLANETEN\par}
PALLAS OLBERSIANA.\par
\medskip\hrule\medskip
Monatliche Correspondenz. Band VI. S.\ 71—96. Juli 1802.\par
\medskip\hrule
\end{center}\medskip
—\quad —\quad —

„Der sehr lehrreiche Aufsatz von Prof.\ \textsc{Wurm} in Ihrem Juniheft (schreibt Dr.\ \textsc{Gauß} am 26. Juni) zeigt leider, dass wir wegen der noch immer unsicheren Masse des Jupiter die Störungsgleichungen der \textit{Ceres} und der \textit{Pallas} noch nicht mit der Zuverlässigkeit werden bestimmen können, die man sich wohl wünschen würde. Dagegen werden die Beobachtungen der \textit{Ceres}, und noch mehr die der \textit{Pallas}, auch den großen Nutzen haben, dass sich jene Masse aus ihnen nach einem oder ein paar Umläufen sehr gut bestimmen lassen wird. Das kann sich dann auf viele andere wichtige Punkte der Astronomie nur günstig auswirken. So reicht in dieser erhabenen Wissenschaft, die der Himmel den Menschen geradezu geschenkt zu haben scheint, damit sie ihre Kräfte entfalten und sich über das Irdische erheben, eine Entdeckung der anderen immer die Hand. Die großen Entdeckungen stehen nie allein, sondern sind oft in ihren Folgen ebenso wichtig wie an sich. Eben jene Unbestimmtheit wird mir einen weiteren Anlass geben, meinem früheren, schon erwähnten Vorsatz treu zu bleiben. Außer der neuen Ellipse für \textit{Ceres}, die wenigstens unter genauester Berücksichtigung der Jupiterstörungen zu berechnen ist, will ich noch eine andere, \textit{reine} Ellipse ohne jede Störung berechnen, die auf genau denselben Beobachtungen beruht. Ich glaube, es müsste doch interessant sein zu sehen, inwieweit sich nächstes Jahr schon eine \textit{sichere} Spur der Einwirkung Jupiters in der Bewegung der \textit{Ceres} zeigen wird. Bis jetzt zeigt sich nämlich noch nicht die geringste Spur, und die Beobachtungen von anderthalb Jahren werden sich ohne Zweifel recht gut und ohne Zwang durch eine \textit{reine Ellipse} darstellen lassen. Dass die Elemente VII im Mai um $40''$ abweichen, braucht einen gar nicht zu wundern, wenn man bedenkt, auf wie dürftigen Beobachtungen sie beruhten. Ich hatte ja nur die Beobachtungen aus Palermo und ein paar von Ihnen, zu denen die Deklinationen fehlten. Aus diesem Grund kann ich meine bisherige Meinung noch nicht aufgeben, dass es die Arbeit ganz unnötig vermehrt hätte, wenn ich schon damals die Störungen hätte berücksichtigen wollen. Ich sehe tatsächlich nicht, dass man dadurch bisher etwas hätte besser machen können. Wohl aber glaube ich, dass vieles schlechter geworden wäre. Ich jedenfalls hätte wohl kaum die Geduld gehabt, eine so große Zahl von Beobachtungen haarscharf mit den Elementen zu vergleichen und eine Ephemeride für eine geraume Zeit zu berechnen, wenn ich jedes Mal noch so viele Gleichungen hätte berechnen müssen. Eine ohne Rücksicht auf die Störungen berechnete Ellipse, die sich genau an die Beobachtungen hält, muss diese Störungen selbst eine geraume Zeit lang mit einschließen. Solange es nur darum geht, die Beobachtungen zu erleichtern, die schon angestellten zu erörtern und ihren jeweiligen Wert zu würdigen, muss uns diese Ellipse weit wichtiger sein als eine \textit{mittlere Ellipse}. Diese kann erst dann interessant werden, wenn die Beobachtungen lange genug gedauert haben, um sie mit solcher Sicherheit bestimmen zu können, dass man mit ihr einige Jahre zurückgehen und in den älteren Sternverzeichnissen nachsuchen kann. Jetzt kann man sich schon nach anderthalb Jahren etwas recht Zuverlässiges davon versprechen. Deshalb werde ich, sobald ich den Schluss der Beobachtungen aus Palermo
\SourceLeaf{223}%
habe, diese Rechnungen mit aller mir möglichen Sorgfalt anstellen. Aus diesen Gründen glaube ich auch, dass es wenig nützen würde, bei \textit{Pallas} schon in diesem Jahr die Störungen zu berücksichtigen, zumal die bisher bei anderen Planeten üblichen Methoden meiner Meinung nach bei diesem Planeten keineswegs die nötige Genauigkeit ergeben würden. Nur wenn die Lichtschwäche des Planeten, \textit{quod Deus avortat}, seine Wiederauffindung im nächsten Jahr vereiteln sollte und man deshalb mit den diesjährigen Beobachtungen zwei Jahre vorausrechnen müsste, wäre es gut, die Einwirkung des Jupiter zu berücksichtigen, wenn auch meines Erachtens auf einem anderen als dem üblichen Weg.“\quad —\quad —\quad —

\Needspace{14\baselineskip}
\begin{center}
\textit{Elemente der Pallas Olbersiana aus 42 Tagen Meridianbeobachtungen auf dem Seeberg und in Greenwich.}\par\medskip
(III.)\par\medskip
{\small\setlength{\tabcolsep}{2pt}
\begin{tabular}{l r}
Epoche 1802, 31. März, Mittag auf dem Seeberg $\ .\ .\ .\ .\ .\ .\ .$&$162^0\ 25'\ 45''9$\\
tägliche mittlere tropische Bewegung $\ .\ .\ .\ .\ .\ .\ .$&$769''547$\\
Logarithmus der großen Halbachse $\ .\ .\ .\ .\ .\ .\ .\ .$&$0.4425664$
\end{tabular}
\par\smallskip
\resizebox{0.98\linewidth}{!}{%
$\left.\begin{array}{l}\text{Sonnenferne}\\\text{Knoten}\quad .\end{array}\right\}\ \text{für die Epoche und gegenüber den Fixsternen ruhend}\left\{\begin{array}{lr}
.\ .&300^0\ 58'\ 47''7\\
.\ .&172^0\ 28'\ 17''9
\end{array}\right.$%
}
\par\smallskip
\begin{tabular}{l r}
Exzentrizität $\ .\ .\ .\ .\ .\ .\ .\ .\ .\ .\ .\ .\ .\ .\ .\ .\ .$&$0.2476402$\\
Neigung $\ .\ .\ .\ .\ .\ .\ .\ .\ .\ .\ .\ .\ .\ .\ .\ .\ .\ .$&$34^0\ 39'\ 10''7$
\end{tabular}}
\end{center}

Es ist doch höchst bemerkenswert, dass die mittlere tropische Bewegung der \textit{Pallas} derjenigen der \textit{Ceres} unbeabsichtigt und ganz von selbst immer näher kommt. Nach den Elementen IV der Ceresbahn (Dezemberheft 1801, S.\ 647 [203 dieses Bandes]) bestimmte Dr.\ \textsc{Gauß} diese mittlere Bewegung zu $770''914$. Sie unterscheidet sich von der oben angegebenen Bewegung der \textit{Pallas} nur um anderthalb Sekunden. Da aber beide Bewegungen noch mit einer gewissen Unsicherheit behaftet sind, kann man jetzt eigentlich noch gar keinen Unterschied angeben.

\Needspace{27\baselineskip}
Dr.\ \textsc{Gauß} verglich nun die folgende Reihe von Beobachtungen mit den obigen Elementen III und fand die ausgezeichnete Übereinstimmung, die die folgende Tabelle zeigt.
\begin{center}
{\small\setlength{\tabcolsep}{4pt}
\begin{tabular}{r r|r|r|r|r|l}
\multicolumn{2}{c|}{1802}&\shortstack{Berechnete\\Rektaszension\\der Pallas}&\shortstack{Berechnete\\Deklination\\der Pallas}&\multicolumn{2}{c|}{Unterschied}&Beobachter.\\
\multicolumn{2}{c|}{}&&&in R.A.&in Dekl.&\\\hline
April&4&$183^0\ 44'\quad5''6$&$13^0\ 54'\ 53''5$&$-1''0$&$+1''5$&von Zach\\
&5&$183\ 34\ 24.4$&$14\ 13\ 20.2$&$+0.$&$-2.7$&von Zach\\
&7&$183\ 15\ 38.5$&$14\ 48\ 59.3$&$0.7$&$-2.8$&von Zach\\
&8&$183\quad6\ 33.9$&$15\quad6\ 12.8$&$-3.9$&$+2.8$&von Zach\\
&10&$182\ 48\ 56.3$&$15\ 39\ 43.5$&$-1.7$&$+1.5$&Dr.\ Burckhardt\\
&12&$182\ 32\ 33.7$&$16\ 11\quad6.4$&$+2.2$&$+8.4$&Méchain\\
&13&$182\ 24\ 46.8$&$16\ 26\quad8.3$&$+2.4$&$+2.3$&Méchain\\
&15&$182\ 10\ 14.1$&$16\ 54\ 33.7$&$-2.4$&$+2.9$&von Zach\\
&15&$182\ 10\quad4.2$&$16\ 54\ 53.3$&$+1.3$&$+1.9$&Méchain\\
&16&$182\quad3\quad9.7$&$17\quad8\ 36.3$&$+1.7$&$-0.5$&Méchain\\
&17&$181\ 56\ 33.7$&$17\ 21\ 52.8$&$+3.4$&$+3.0$&Méchain\\
&18&$181\ 50\ 25.0$&$17\ 34\ 25.3$&$-5.6$&$\ .\ .\ .$&von Zach\\
&19&$181\ 44\ 26.4$&$17\ 46\ 50.0$&$+1.1$&$-4.4$&von Zach\\
&23&$181\ 23\ 41.9$&$18\ 32\ 27.4$&$-0.1$&$-0.6$&Dr.\ Maskelyne\\
&24&$181\ 19\ 32.4$&$18\ 42\ 25.2$&$-5.4$&$\ .\ .\ .$&von Zach\\
&25&$181\ 15\ 35.3$&$18\ 52\ 15.9$&$+3.1$&$\ .\ .\ .$&von Zach\\
&25&$181\ 15\ 28.6$&$18\ 52\ 33.1$&$+0.3$&$+3.4$&Dr.\ Maskelyne\\
&26&$181\ 11\ 59.4$&$19\quad1\ 41.6$&$+2.4$&$-7.8$&von Zach\\
&26&$181\ 11\ 53.3$&$19\quad1\ 58.1$&$+3.1$&$-4.5$&Dr.\ Maskelyne\\
&27&$181\quad8\ 44.8$&$19\ 10\ 42.7$&$+0.8$&$-4.1$&von Zach
\end{tabular}}
\end{center}
\expandafter\def\csname HistoricalNoteText20\endcsname{Bei der \textit{Greenwicher} Beobachtung der R.A. vom 7. Mai scheint irgendein Versehen vorgekommen zu sein, weil sie um etwa $39''$ zu klein zu sein scheint.}
\SourceLeaf{224}%
\begin{center}
{\small\setlength{\tabcolsep}{4pt}
\begin{tabular}{r r|r|r|r|r|l}
\multicolumn{2}{c|}{1802}&\shortstack{Berechnete\\Rektaszension\\der Pallas}&\shortstack{Berechnete\\Deklination\\der Pallas}&\multicolumn{2}{c|}{Unterschied}&Beobachter\\
\multicolumn{2}{c|}{}&&&in R.A.&in Dekl.&\\\hline
April&29&$181^0\quad3'\ 20''9$&$19^0\ 27'\ 31''5$&$+4''3$&$-12''2$&von Zach\\
&30&$181\quad1\ 12.1$&$19\ 35\ 19.6$&$+2.2$&$-13.9$&von Zach\\
Mai&1&$180\ 59\ 24.6$&$19\ 42\ 44.8$&$+2.1$&$+4.0$&von Zach\\
&1&$180\ 59\ 21.7$&$19\ 42\ 57.7$&$+2.8$&$+3.3$&Dr.\ Maskelyne\\
&2&$180\ 57\ 59.7$&$19\ 49\ 46.5$&$+7.8$&$-5.5$&von Zach\\
&2&$180\ 57\ 57.5$&$19\ 49\ 58.7$&$+0.2$&$+2.0$&Dr.\ Maskelyne\\
&3&$180\ 56\ 56.9$&$19\ 56\ 25.9$&$+1.3$&$-21.8$&von Zach\\
&4&$180\ 56\ 15.0$&$20\quad2\ 53.0$&$+2.1$&$-0.3$&Dr.\ Maskelyne\\
&5&$180\ 55\ 56.9$&$20\quad8\ 36.2$&$+2.3$&$-2.5$&von Zach\\
&6&$180\ 56\quad0.3$&$20\ 14\quad8.8$&$+6.2$&$\ .\ .\ .$&von Zach\\
&7&$180\ 56\ 24.7$&$20\ 19\ 20.4$&$+0.9$&$-14.9$&von Zach\\
&7&$180\ 56\ 25.8$&$20\ 19\ 29.4$&$+39.3$&$-0.5$&Dr.\ Maskelyne\\
&8&$180\ 57\ 13.0$&$20\ 24\ 10.8$&$+3.0$&$-9{,}3$&von Zach\\
&9&$180\ 58\ 23.9$&$20\ 28\ 49.3$&$+2.3$&$+6.8$&Dr.\ Maskelyne\\
&11&$181\quad1\ 44.6$&$20\ 36\ 41.9$&$+2.9$&$-31.3$&von Zach\\
&11&$181\quad1\ 48.3$&$20\ 36\ 48.5$&$+3.4$&$+0.3$&Dr.\ Maskelyne\\
&13&$181\quad6\ 37.6$&$20\ 43\ 31.4$&$+4.6$&$-1.7$&Dr.\ Maskelyne\\
&14&$181\quad9\ 33.4$&$20\ 46\ 25.2$&$+4.0$&$-1.5$&Dr.\ Maskelyne\\
&16&$181\ 16\ 26.5$&$20\ 51\ 19.3$&$-1.7$&$+0.5$&Dr.\ Maskelyne
\end{tabular}}
\end{center}

Dieser Vergleich zeigt, dass man nach dieser ganzen Beobachtungsreihe nur noch wenig an den Elementen ändern könnte, zumal die letzten Greenwicher Deklinationen im Mai (die einen größeren Einfluss auf die Elemente haben als die Rektaszensionen) noch sehr gut miteinander übereinstimmen. Nur die berechneten Rektaszensionen\HistoricalNote{20} müssten gegen Ende noch um etwa $4''$ kleiner gemacht werden. Die Elemente wären davon aber nur ganz unerheblich betroffen. In einem Schreiben vom 26. Juni teilt uns Dr.\ \textsc{Gauß} mit, er habe auch die Beobachtungen dieses Planeten durch \textsc{Olbers} vom 19., 20. und 21. Juni noch mit seinen Elementen III verglichen. Sie stimmten noch so schön überein, dass er danach mit Gewissheit gar nichts ändern könne. Man wird es demnach bei diesen Bahnelementen III belassen können, bis die Beobachtungen ganz abgeschlossen sind. Damit berechnete Dr.\ \textsc{Gauß} also die folgende
\Needspace{15\baselineskip}
\begin{center}
\textit{Ephemeride für die Pallas Olbersiana für Mitternacht auf dem Seeberg.}\par\medskip
{\small\setlength{\tabcolsep}{3pt}
\begin{tabular}{r r|r|r|r|r}
\multicolumn{2}{c|}{1802}&\shortstack{R.A.\\der Pallas}&\shortstack{Dekl.\\der Pallas}&\shortstack{Abstand\\von d.\ $\EarthSymbol$}&\shortstack{Abstand\\von d.\ $\odot$}\\\hline
Mai&18&$181^0\ 25'$&$20^0\ 55'$&$1.888$&$2.477$\\
&21&$181\ 40$&$20\ 59$&$1.930$&$2.484$\\
&24&$181\ 58$&$21\quad0$&$1.973$&$2.492$\\
&27&$182\ 19$&$20\ 59$&$2.016$&$2.500$\\
&30&$182\ 42$&$20\ 56$&$2.060$&$2.508$\\
Juni&2&$183\quad8$&$20\ 51$&$2.104$&$2.516$\\
&5&$183\ 36$&$20\ 44$&$2.148$&$2.524$\\
&8&$184\quad6\frac{1}{2}$&$20\ 36$&$2.192$&$2.532$
\end{tabular}
\par\medskip
\begin{tabular}{r r|r|r|r|r}
\multicolumn{2}{c|}{1802}&\shortstack{R.A.\\der Pallas}&\shortstack{Dekl.\\der Pallas}&\shortstack{Abstand\\von d.\ $\EarthSymbol$}&\shortstack{Abstand\\von d.\ $\odot$}\\\hline
Juni&11&$184^0\ 39'$&$20^0\ 26\frac{1}{2}'$&$2.237$&$2{,}539$\\
&14&$185\ 14$&$20\ 15$&$2.281$&$2.547$\\
&17&$185\ 51$&$20\quad3$&$2.326$&$2.555$\\
&20&$186\ 30$&$19\ 50$&$2.370$&$2.563$\\
&23&$187\ 10$&$19\ 35$&$2.415$&$2.571$\\
&26&$187\ 52\frac{1}{2}$&$19\ 19\frac{1}{2}$&$2.459$&$2.579$\\
&29&$188\ 36$&$19\quad3$&$2.504$&$2.587$
\end{tabular}}
\end{center}

Für den Anfang des kommenden Jahres berechnete Dr.\ \textsc{Gauß} den ungefähren Ort der \textit{Pallas} wie folgt. Am 31. Dez.\ um $12^{\mathrm{u}}$ Mitternacht auf dem Seeberg, Länge $253^0\ 33'$, Breite $26^0\ 10'$, nördl.\ Abstand von der Erde $3.779$, von der Sonne $3.043$. Es bleibt aber sehr ungewiss, ob wir diesen Planeten wegen seiner Lichtschwäche
\SourceLeaf{225}%
in den nächsten beiden Jahren wiedersehen werden. Dr.\ \textsc{Olbers} hat die beiden folgenden Oppositionen der \textit{Pallas} nach den Elementen II von \textsc{Gauß} berechnet. Als er sie am 28. März entdeckte, war ihr Abstand von der Erde nur $=1.38$. Ende Mai war er $=2.00$ und am 15. Juni $=2.31$. Im kommenden Jahr 1803 gerät \textit{Pallas} gegen Ende Juni wieder in Opposition zur Sonne. Sie wird über dem Cerberus bei $9^{\mathrm{z}}\ 8^0$ Länge und $48^0$ nördlicher Breite stehen, aber ihr Abstand von der Erde wird $2.61$ betragen. Dr.\ \textsc{Olbers} glaubt, dass man sie bei diesem Abstand in der hellen nördlichen Dämmerung kaum finden dürfte. Etwas günstiger ist die Opposition von 1804, die gegen Ende August bei $11^{\mathrm{z}}\ 3^0$ und $19^0$ nördlicher Breite eintritt, also im Kopf des Pegasus. Zwar findet Dr.\ \textsc{Olbers} für den Abstand von der Erde noch $2.40$, doch die Nächte sind dann schon vollkommen dunkel. Aber \textit{ganz gewiss}, glaubt Dr.\ \textsc{Olbers}, werden wir \textit{Pallas} im Jahr 1805 wiedersehen, da sie bei der Opposition schon in den sonnennäheren Teil ihrer Bahn herabgekommen ist und damit auch der Erde viel näher sein wird. Hoffentlich halten diese Überlegungen unterdessen keinen Astronomen davon ab, \textit{Pallas} nach Kräften aufzuspüren, sobald eine Suche möglich ist. Denn große, lichtstarke Fernrohre könnten doch ausreichen, diesen Planeten noch zu entdecken, selbst wenn er nur als Stern zwölfter Größe erscheint. Freilich wird es etwas schwieriger sein, ihn zu finden. Doch selbst wenn 1803 und 1804 alle Versuche fehlschlagen sollten, hofft Dr.\ \textsc{Gauß}, dass seine Bahnelemente für 1805 immer noch als Leitfaden ausreichen werden, um \textit{Pallas} dennoch wiederzufinden (freilich mit etwas mehr Mühe als bei \textit{Ceres}).\quad —\quad —\quad —

\Needspace{30\baselineskip}
Was die geringe Größe dieser beiden neuen planetarischen Himmelskörper angeht, so ist es tatsächlich schwer zu begreifen, wie man einen so kleinen scheinbaren Durchmesser genau messen kann. \textsc{Herschel} und \textsc{Schröter}, beide geübte Beobachter und mit den größten und besten optischen Instrumenten ausgestattet, haben solche Messungen versucht. Schröter fand \textit{Ceres} 18-mal und \textit{Pallas} 40-mal so groß wie Herschel. Um die Schwierigkeit solcher Messungen zu zeigen, hat Dr.\ \textsc{Gauß} nach seinen neuesten Bahnelementen die folgenden Berechnungen der scheinbaren Durchmesser dieser Himmelskörper angestellt.
\begin{center}
\textit{Durchmesser der Ceres.}\par\medskip
{\small\setlength{\tabcolsep}{6pt}
\begin{tabular}{r r|r|c}
\multicolumn{2}{c|}{1802}&\shortstack{Abstand\\von der\\$\EarthSymbol$}&\shortstack{Scheinbarer Durchmesser, wenn man\\den wahren nach Dr.\ \textsc{Herschel}\\mit $=35$ deutschen Meilen ansetzt.}\\\hline
März&28&$1.621$&$0''22$
\end{tabular}}
\par\bigskip
\textit{Durchmesser der Pallas.}\par\medskip
{\small\setlength{\tabcolsep}{6pt}
\begin{tabular}{r r|r|c}
\multicolumn{2}{c|}{1802}&\shortstack{Abstand\\von der\\$\EarthSymbol$}&\shortstack{Scheinbarer Durchmesser, wenn man\\Dr.\ \textsc{Herschel}s Ergebnis $=15$ deutsche\\Meilen für richtig hält.}\\\hline
April&4&$1.407$&$0''11$\\
—&15&$1.386$&$0.10$\\
Juni&23&$2.415$&$0.06$
\end{tabular}}
\end{center}

Ob man so kleine Durchmesser wirklich noch messen oder überhaupt das reflektierte Licht von so kleinen Oberflächen noch sehen kann, überlassen wir anderen zur Beurteilung.\quad —\quad —\quad $=$
\begin{flushright}[\textsc{von Zach}.]\end{flushright}
